当然, 若函数属于 $W^{1,p}(\Omega)$, 其差商就受到偏导数的某种界的控制, 这毫不奇怪. 以下结果更为实质: 若开始只知道 $u$ 属于 $L^p(\Omega)$, 并能得到与 $h$ 无关的差商界, 则可推出 $u$ 属于空间 $W^{1,p}(\Omega)$.

\begin{Lemma}
\label{lem:uniform-difference-quotient-converse}
设 $u\in L^p(\Omega)$, $1<p\leq\infty$, 并假设存在常数 $\widetilde C$, 使对任意 $\Omega^{\prime}\subset\subset\Omega$ 和 $h<\operatorname{dist}(\Omega^{\prime},\partial\Omega)$, 有 $D_i^hu\in L^p(\Omega^{\prime})$, 且
\begin{equation}\label{eq:uniform-difference-quotient-converse-hypothesis}
\|D_i^hu\|_{L^p(\Omega^{\prime})}\leq\widetilde C.
\end{equation}
则 $u_{x_i}$ (先验地作为分布已有良好定义) 实际上属于 $L^p(\Omega)$, 并满足
\begin{equation}\label{eq:uniform-difference-quotient-recovered-derivative-bound}
\|u_{x_i}\|_{L^p(\Omega)}\leq\widetilde C.
\end{equation}

\end{Lemma}

\begingroup
\renewcommand{\endProof}{\par\nobreak\hfill$\blacksquare$\par\vspace{0.1cm}}
\begin{Proof}
回顾在 Section~\ref{subsec:boundary-value-problems-adjoints} 末尾已经说明, 函数的分布导数可由差商的极限得到. 在当前问题的记号下, 有
\begin{equation}\label{eq:difference-quotient-distributional-derivative-limit}
(u_{x_i},\phi)=\lim_{h\to0}(D_i^hu,\phi)
\end{equation}
对每个 $\phi\in\mathcal{D}(\Omega)$ 成立.

现在注意, 存在函数 $v\in L^p(\Omega)$, 满足 $\|v\|_{L^p(\Omega)}\leq\widetilde C$, 以及序列 $h_m\to0$, 使对每个 $\phi\in\mathcal{D}(\Omega)$ 都有
\begin{equation}\label{eq:difference-quotient-weak-limit-function}
\lim_{h_m\to0}\int_\Omega\phi D_i^{h_m}u\,\mathrm{d}\bm{x}=\int_\Omega\phi v\,\mathrm{d}\bm{x}.
\end{equation}
这来自对每个 $\Omega^{\prime}\subset\subset\Omega$, 有界集 $D_i^hu$ 在 $L^p(\Omega^{\prime})$ 中的弱紧性 (参见 Theorem~\ref{thm:alaoglu-weak-star-subsequence}), 以及以下事实: 可以用可数个闭子集 $\Omega_j\subset\subset\Omega$ 覆盖 $\Omega$, 使每个 $\Omega^{\prime}\subset\subset\Omega$ 至多与其中有限个 $\Omega_j$ 相交.

因此,
\begin{equation}\label{eq:difference-quotient-recovered-distributional-derivative}
(u_{x_i},\phi)=\int_\Omega\phi v\,\mathrm{d}\bm{x};
\end{equation}
即 $u$ 的分布导数由函数 $v\in L^p(\Omega)$ (唯一地) 给出.
\end{Proof}
\endgroup

还需要建立使用差分算子的几个重要工具, 即微分学中乘积法则和分部积分的对应形式.

\begin{Lemma}
\label{lem:difference-quotient-product-and-parts}
假设 $\Omega^{\prime}\subset\subset\Omega$ 且 $h<\operatorname{dist}(\Omega^{\prime},\partial\Omega)$. 则对任意 $u\in L^p(\Omega)$, $1<p<\infty$, 以及任意满足 $\operatorname{supp}v\subseteq\Omega^{\prime}$ 的 $v\in C_0(\Omega)$, 有
\begin{equation}\label{eq:difference-quotient-shifted-product-rule}
D_i^h(uv)(\bm{x})=u(\bm{x})D_i^hv(\bm{x})+D_i^hu(\bm{x})v(\bm{x}+h\bm{e}_i)
\end{equation}
和
\begin{equation}\label{eq:difference-quotient-integration-by-parts}
\int_\Omega D_i^hu(\bm{x})v(\bm{x})\,\mathrm{d}\bm{x}=-\int_\Omega u(\bm{x})D_i^{-h}v(\bm{x})\,\mathrm{d}\bm{x}.
\end{equation}

\end{Lemma}

证明留给读者.
